% Generated draft: requires mathematical and source review.
\documentclass[11pt]{article}
\usepackage[margin=25mm]{geometry}
\usepackage{fontspec}
\setmainfont{Latin Modern Roman}
\usepackage{amsmath,amssymb,mathtools,unicode-math}
\setmathfont{Latin Modern Math}
\usepackage{enumitem,needspace,xurl,hyperref}
\setlength{\parindent}{0pt}
\setlength{\parskip}{5pt}
\setlength{\emergencystretch}{3em}
\newcommand{\N}{\mathbb{N}}
\newcommand{\Z}{\mathbb{Z}}
\newcommand{\Q}{\mathbb{Q}}
\newcommand{\R}{\mathbb{R}}
\newcommand{\C}{\mathbb{C}}
\newcommand{\F}{\mathbb{F}}
\newcommand{\A}{\mathbb{A}}
\newcommand{\PP}{\mathbb{P}}
\newcommand{\E}{\mathbb{E}}
\newcommand{\eps}{\varepsilon}
\newcommand{\dd}{\,\mathrm{d}}
\begin{document}
\section*{UnsolvedMath ID: 30006285}
\textbf{Source problem code:} catalogue entry
\subsection{Definitions and mathematical statement}
Write \(\operatorname{CLE}_{\kappa}\) for the conformal loop ensemble with parameter \(\kappa\). The supplied record does not specify the range of \(\kappa\), the domain or boundary conditions of the ensemble, or the construction and underlying space of the object called the \(\operatorname{CLE}_{\kappa}\) metric.

For a connected graph \(G=(V,E)\) with unit edge lengths, its intrinsic, or chemical-distance, metric is
\[
 d_G(x,y)=\inf\bigl\{n\in\Z_{\geq 0}:\text{there exist }v_0=x,\ldots,v_n=y\text{ with }\{v_{i-1},v_i\}\in E\text{ for }1\leq i\leq n\bigr\}.
\]
This defines distance within a connected graph or cluster. The source does not specify which graphs or clusters supply the metric spaces in the conjecture, whether other edge-length conventions are allowed, or how disconnected components are treated.

\textbf{Conjecture as supplied.} For every \(\kappa\), the \(\operatorname{CLE}_{\kappa}\) metric describes the scaling limit of the intrinsic metric of every discrete model whose interfaces converge to \(\operatorname{CLE}_{\kappa}\) in the scaling limit. In the original wording, the hypothesis is that the discrete model converges to \(\operatorname{CLE}_{\kappa}\); the curated statement makes interface convergence explicit.

The proposed conclusion can only be expressed schematically from this record: if \(G_\delta\) denotes the relevant discrete graph or cluster at mesh \(\delta\), and \(D_\kappa\) denotes the intended continuum metric, then suitable distance normalizations \(a_\delta>0\) should give
\[
 (V(G_\delta),a_\delta d_{G_\delta})\longrightarrow (X_\kappa,D_\kappa)
 \qquad\text{as }\delta\downarrow 0.
\]
Here \(X_\kappa\), \(D_\kappa\), the normalization, and the meaning of convergence are placeholders for specifications absent from the record. The source does not say whether normalizations are deterministic, how they depend on the model, whether the limit is unique up to a multiplicative constant, or whether metric convergence must hold jointly with interface convergence. Consequently this display records the intended universality claim, rather than a fully specified convergence assertion. The source's universal quantifiers are retained; their parameter range and admissible model class remain unresolved.
\subsection{Short English statement}
The conjecture asks whether every discrete model with a \(\operatorname{CLE}_{\kappa}\) interface limit also has the corresponding continuum limit for suitably rescaled chemical distances. The parameter range, admissible models, metric space, normalization, and convergence notion require specification. Supplied dated curation labels the problem partially solved and reports a critical planar percolation result at \(\kappa=6\); that assessment has not been independently verified.
\subsection{Sources}
\begin{itemize}
\item Local UnsolvedMath registry, permanent ID \(30006285\), title ``Universal CLE Metrics as Chemical-Distance Limits,'' exact source problem code \texttt{catalogue entry}. The supplied record gives proposed year 2025 and no separately identified proposer.
\item \emph{Statistical Physics Out of Equilibrium: Quantitative Results and Universality} (2025), Oberwolfach Reports, DOI 10.4171/owr/2025/26, pp.\ 1333--1388. Supplied source report and literature-source link: \url{https://doi.org/10.4171/owr/2025/26}.
\item The original excerpt associates the conjecture with Jason Miller's contribution, ``The scaling limit of simple random walk and the intrinsic metric on 2D critical percolation clusters,'' described as joint work with Valeria Ambrosio, Irina Dankovic, Maarten Markering, and Yizheng Yuan. This preserves the supplied contribution attribution without identifying a separate conjecture proposer.
\item Supplied literature source: Dankovic--Markering--Miller--Yuan, ``planar critical percolation,'' \url{https://arxiv.org/abs/2604.14122}. The record supplies this descriptive label, rather than a complete bibliographic citation.
\item Literature assessment checked at \texttt{2026-08-22}: the supplied curation labels the problem \texttt{partially\_solved}, states that the universal \(\operatorname{CLE}_{\kappa}\) metric conjecture remains open, and reports that the cited 2026 work proves intrinsic-metric convergence for critical planar percolation, identified there as the \(\kappa=6\) case. This is supplied dated curation derived from the OWR record, not an independently verified assessment of the literature or of the cited theorem's hypotheses and scope.
\end{itemize}
\end{document}
